\section{LANDASAN TEORI}
\textbf{Headline Incongruity Detection (HID)}. HID didefinisikan sebagai persoalan menilai apakah informasi utama pada judul berita secara faktual didukung oleh isi berita yang menyertainya. Persoalan ini berbeda dari deteksi \textit{clickbait}, yang berfokus pada judul provokatif atau sensasional terlepas dari kebenarannya \cite{chesney17}, karena HID menuntut penilaian terhadap hubungan referensial antara dua unit teks, mendekati kerangka \textit{textual entailment} dan identifikasi parafrasa. Riset awal memperlakukan HID sebagai persoalan klasifikasi berbasis fitur buatan tangan, sementara riset lanjutan bergeser ke representasi semantik yang dipelajari dari data \cite{yoon21, bourgonje17}. Pada konteks bahasa Indonesia, kajian menunjukkan pola kegagalan yang serupa dengan bahasa Inggris, bahkan dengan tantangan tambahan berupa morfologi aglutinatif dan keterbatasan sumber daya berlabel \cite{william20, surianto26}. Kesamaan lintas bahasa ini menegaskan bahwa akar persoalan HID bukan pada bahasa tertentu, melainkan pada keterbatasan mendasar ukuran kemiripan permukaan dalam menangkap perubahan makna.

\textbf{\textit{Vector Space Model} \& \textit{Lexical Similarity}}.
Model ruang vektor dengan pembobotan TF-IDF \cite{salton88} dan ukuran himpunan seperti Jaccard mengasumsikan teks yang berbagi kosakata membicarakan hal yang sama, dan terbukti menjadi \textit{baseline} kuat pada deteksi sikap berita \cite{riedel17}, tetapi rapuh terhadap pertukaran entitas atau angka yang mempertahankan kata yang sama. Hipotesis distribusional, bahwa kata pada konteks serupa memiliki makna berdekatan \cite{harris54}, mendasari \textit{Word2Vec} yang mempelajari vektor kata dari korpus tanpa label eksternal \cite{mikolov13}, sementara \textit{pointwise mutual information} mengukur kekuatan asosiasi dua kata dari statistik kemunculan bersama \cite{church90}. Kedua kerangka ini memungkinkan estimasi kedekatan makna sepenuhnya dari data latih tanpa bobot pralatih.

\textbf{Penyelarasan Urutan dan Pemeringkatan Relatif}. Selain kosakata, urutan informasi membawa makna struktural. Penyelarasan rangkaian yang berakar pada pencocokan biologi komputasional \cite{needleman70} dan korelasi peringkat Kendall tau \cite{kendall38} menyediakan kerangka untuk mengukur konsistensi urutan, relevan karena berita konvensional mengikuti gaya piramida terbalik sehingga penyimpangan pola monoton menjadi sinyal struktural. Teori \textit{learning to rank} \cite{liu09}, yang berakar pada hukum penilaian komparatif Thurstone \cite{thurstone27}, menunjukkan bahwa perbandingan relatif antar kandidat sering lebih diskriminatif daripada penilaian mutlak satu kandidat secara terisolasi.

\textbf{Encoder Neural, Atensi, dan Augmentasi Data}. \textit{Gated recurrent unit} memodelkan ketergantungan sekuensial secara hemat parameter \cite{cho14}, sedangkan mekanisme atensi membentuk penyelarasan lunak dinamis antar-elemen dua rangkaian \cite{bahdanau14}. \textit{Enhanced Sequential Inference Model} menggabungkan keduanya melalui atensi silang sebelum komposisi \cite{chen17}, merepresentasikan paradigma \textit{cross-encoder} yang menangkap interaksi berbutir halus, berbeda dari \textit{bi-encoder} yang mengorbankan kedalaman interaksi demi efisiensi. Untuk memperkaya data latih pada kelas minoritas, pembangkitan contoh negatif melalui perturbasi terkendali \cite{vincent08} dapat diarahkan pada tipologi kesalahan nyata seperti pertukaran entitas, negasi, dan permutasi urutan, tanpa bergantung pada model generatif pralatih.

